\documentclass{article}
\usepackage[T1]{fontenc}
\usepackage{lmodern}
\usepackage{graphicx}
\usepackage{amsmath,amssymb}
\usepackage[a4paper,margin=2cm]{geometry}
\usepackage[absolute,overlay]{textpos}
\setlength{\TPHorizModule}{1mm}
\setlength{\TPVertModule}{1mm}
\usepackage[indonesian]{babel}
\usepackage{hyperref}
\setlength{\emergencystretch}{2em}
% unit: Halaman 32

\begin{document}

% === Halaman 32 (python/native) ===

32

Kelebihan Mode CBC


 Blok-blok plainteks yang sama tidak selalu menghasilkan blok-blok 

cipherteks yang sama

 Oleh karena blok-blok plainteks yang sama tidak menghasilkan 

blok-blok cipherteks yang sama, maka kriptanalisis menjadi lebih 

sulit. 

 Inilah alasan utama penggunaan mode CBC.

\end{document}
